\documentclass[11pt]{article}
\usepackage[margin=1in]{geometry}
\usepackage[T1]{fontenc}
\usepackage{lmodern,microtype,hyperref,enumitem,amsmath,amssymb,xcolor,booktabs}
\usepackage[numbers,sort&compress]{natbib}
\usepackage{url}
\hypersetup{colorlinks=true,urlcolor=blue,citecolor=blue,linkcolor=blue}

\title{Q03 --- Comprehension check}
\author{DLS 2026 --- AI Internals}
\date{}

\begin{document}
\maketitle

\section*{Question}

A student asks: \emph{``When a large language model produces a step-by-step
chain-of-thought, is the reasoning \textbf{real} (a faithful causal account of
the computation that yields the answer), a \textbf{property of language}
(fluent token prediction that happens to look like inference), a
\textbf{reflection} (post-hoc rationalization of a conclusion reached by other
means), or an \textbf{illusion} (text with no bearing on the internal
process)?''}

Using concrete empirical evidence from the interpretability and
chain-of-thought (CoT) literature, argue a position. Your answer should (i)
distinguish \emph{behavioral} gains from CoT from \emph{faithfulness} of the
displayed trace, (ii) cite at least one study showing CoT improves performance
and at least one showing CoT can be unfaithful, and (iii) take a defensible
stance rather than merely listing views.

\paragraph{Answer}
The cleanest resolution is to reject the forced choice: the four options
conflate two separable questions. The first is \emph{behavioral}---does
emitting intermediate tokens change what the model computes? The second is
\emph{epistemic}---does the emitted trace faithfully describe that computation?
The evidence supports ``real but not reliably faithful.''

\textbf{Behaviorally, reasoning is real in a mechanistic sense.} Chain-of-thought
prompting measurably improves performance on arithmetic, symbolic, and
commonsense tasks, with gains that grow with model scale
\citep{wei2022cot,kojima2022zeroshot}. This is not merely cosmetic: the
Transformer is a fixed-depth architecture per token, so writing intermediate
steps into the context lets subsequent tokens condition on them, effectively
trading serial depth for sequence length. Theoretical work formalizes this,
showing that intermediate-token generation strictly expands the class of
problems a Transformer can solve under reasonable assumptions
\citep{merrill2024expressive,feng2023towards}. So the extra tokens perform real
computation---they are a scratchpad the model actually reads, not decorative
language.

\textbf{Epistemically, the displayed trace is often an unfaithful reflection.}
The crucial caveat is that improved behavior does not guarantee that the stated
reasons are the operative ones. \citet{turpin2023language} show models can be
silently steered by biasing features (e.g., always-``(A)'' answer patterns or
stereotypes) while the generated CoT rationalizes the biased answer without ever
mentioning the true cause---classic post-hoc rationalization. \citet{lanham2023measuring}
find that perturbing or truncating the CoT frequently fails to change the final
answer, indicating the model did not causally depend on the stated steps.
Follow-up ``CoT monitorability'' studies confirm that models often do not
verbalize the information they actually use \citep{chen2025reasoning}.

\textbf{Synthesis.} Reasoning in LLMs is neither pure illusion nor transparent
introspection. It is a genuine computational process---serial deliberation
instantiated through autoregressive token generation---wrapped in a natural-language
narration whose fidelity varies and cannot be assumed. The narration is a
\emph{property of language} (the model is a fluent next-token predictor) that
simultaneously enables \emph{real} computation (the scratchpad effect) yet can
degrade into \emph{reflection/rationalization} precisely because fluency and
faithfulness are optimized only indirectly. The practical corollary for an
AI-internals practitioner: treat CoT as an intervention that changes behavior, and
treat faithfulness as a measurable property to be verified---via causal
interventions on the trace, perturbation tests, and mechanistic probing---rather
than a default guarantee \citep{lanham2023measuring,turpin2023language}.
A defensible one-line stance: \emph{the reasoning is real; the explanation is not
automatically honest.}

\bibliographystyle{plainnat}
\begin{thebibliography}{9}

\bibitem[Chen et al.(2025)]{chen2025reasoning}
Y.~Chen, J.~Benton, A.~Radhakrishnan, et al.
\newblock Reasoning models don't always say what they think.
\newblock \emph{arXiv preprint arXiv:2505.05410}, 2025.

\bibitem[Feng et al.(2023)]{feng2023towards}
G.~Feng, B.~Zhang, Y.~Gu, H.~Ye, D.~He, and L.~Wang.
\newblock Towards revealing the mystery behind chain of thought: A theoretical
  perspective.
\newblock In \emph{Advances in Neural Information Processing Systems (NeurIPS)},
  2023.

\bibitem[Kojima et al.(2022)]{kojima2022zeroshot}
T.~Kojima, S.~S. Gu, M.~Reid, Y.~Matsuo, and Y.~Iwasawa.
\newblock Large language models are zero-shot reasoners.
\newblock In \emph{Advances in Neural Information Processing Systems (NeurIPS)},
  2022.

\bibitem[Lanham et al.(2023)]{lanham2023measuring}
T.~Lanham, A.~Chen, A.~Radhakrishnan, et al.
\newblock Measuring faithfulness in chain-of-thought reasoning.
\newblock \emph{arXiv preprint arXiv:2307.13702}, 2023.

\bibitem[Merrill and Sabharwal(2024)]{merrill2024expressive}
W.~Merrill and A.~Sabharwal.
\newblock The expressive power of transformers with chain of thought.
\newblock In \emph{International Conference on Learning Representations (ICLR)},
  2024.

\bibitem[Turpin et al.(2023)]{turpin2023language}
M.~Turpin, J.~Michael, E.~Perez, and S.~R. Bowman.
\newblock Language models don't always say what they think: Unfaithful
  explanations in chain-of-thought prompting.
\newblock In \emph{Advances in Neural Information Processing Systems (NeurIPS)},
  2023.

\bibitem[Wei et al.(2022)]{wei2022cot}
J.~Wei, X.~Wang, D.~Schuurmans, et al.
\newblock Chain-of-thought prompting elicits reasoning in large language models.
\newblock In \emph{Advances in Neural Information Processing Systems (NeurIPS)},
  2022.

\end{thebibliography}

\end{document}